\documentclass{beamer}
\usetheme{Madrid}
\usecolortheme{default}
\usepackage{graphicx, hyperref, xcolor}
\usepackage[warn]{mathtext}
\usepackage[T2A]{fontenc}
\usepackage[utf8]{inputenc}
\usepackage[english,russian]{babel}
\usepackage{indentfirst}
\usepackage{subcaption}
\usepackage{float}
\usepackage{amsmath}
\usepackage{mathrsfs}
\usepackage{amsfonts}
\usepackage{amssymb}
\usepackage{booktabs}
\usepackage{tikz}
\usetikzlibrary{arrows.meta, positioning}
\usepackage{graphicx}

\setbeamertemplate{navigation symbols}{}
\setbeamertemplate{footline}[frame number]

\title{Анализ дактилоскопических изображений на основе непрерывной скелетизации}
\author{Семенов Артём Григорьевич}
\institute{ВМК МГУ \\[1em] \textit{Научный руководитель}: д.т.н., Местецкий Леонид Моисеевич}
\date{2025}


\begin{document}

\frame{\titlepage}

\begin{frame}{Введение}
    \begin{itemize}
        \item Задачи, связанные с идентификацией дактилоскопического изображения, или отпечатка пальца, могут возникать в таких ситуациях, как предоставление доступа к закрытым местам, установление личности, определение хозяина отпечатка.
        \item У имеющихся алгоритмов есть проблемы: неустойчивость и
низкая скорость работы.
        \item Предлагается рассмотреть подход, основанный на анализе скелетного графа изображеня.
    \end{itemize}
\end{frame}

% --------------------
\begin{frame}{Данные}
    \begin{columns}[T]
      \begin{column}{0.55\textwidth}
      В качестве даннных использовался датасет $DB2 \_ B$ из соревнования по распознаванию отпечатков пальцев $FVC2002$. Он состоит из 80 изображений, снятых с 10 разных пальцев, ровно 8 изображений соответствуют одному. Сами изображения имеют размер 560 пикселей в высоту и 296 пикселей в ширину.
      \end{column}
      \begin{column}{0.4\textwidth}
        \begin{figure}[H]
            \centering
            \includegraphics[scale=0.25]{new_image_0.png}
            \caption{Пример изображения}
            \label{fig:orange}
         \end{figure}
      \end{column}
    \end{columns}
\end{frame}

% --------------------
\begin{frame}{Постановка задачи}
    

    \vspace{0.5cm}

    \textbf{Дано:}
    Выборки $X_{train} = \{(x_{i}, y_{i})\}_{i = 1}^{\ell}, X_{test} = \{(x_{i}^{'})\}_{i = 1}^{\ell^{'}}; x_{i}, x_{i}^{'} \in X$, где $  X = \{x: x \in \mathbb{N}_{0}^{n \times m}, x \leq 255\}, y_{i} \in Y, $ где $Y$ — конечное множество классов. $\exists f(x): X \to Y$, $\forall y \in Y \exists x_{1} \in X,  x_{2} \in X_{test}: (x_{1}, y) \in X_{train}, f(x_{2}) = y $,
    T — время вычислений функий, зависящих от тестовой выборки, S — некоторая функция штрафа по времени.
    
    
    \vspace{0.3cm}
    
    \textbf{Найти:}
    Функию $g(x)$, приближающую $f(x)$, не требующую слишком много времени на вычисление.
    
    \vspace{0.3cm}
    
    \textbf{Критерий:}
    Средняя точность классификации объектов, со штрафом по времени: $\frac{\sum_{x \in {X_{test}}}(g(x) == f(x))}{|X_{test}|} + S(T)$
\end{frame}


% --------------------
\begin{frame}{Подходы к решению}
    Наиболее распространные подходы
    \begin{itemize}
        \item подходы, ориентированные на использовании минюций
        \item неиросетевые
        \item сравнение отпечатков как целостного образа
        
    \end{itemize}
    В настоящее время, для задачи идентификации самым распространенным является подход, основанный на минюциях. Далее, в эксперименте будет протестирована одна из его реализаций.
    

\end{frame}



% --------------------
\begin{frame}{Эксперимент}
    Весь алгоритм состоит из нескольких шагов:
    \begin{enumerate}
        \item Предобработка изображения. К изображению применяются классические методы обработки, в результате этого папиллярные линии становятся более отчетливыми.  
        \item Выделяются ключевые точки. Множество ключевых точек содержит в себе множество минюций, что используется далее.
        \item Циклом по элементов из объектов датасета, ключевые точки идентифицируемого изображения сопоставляются ключевым точкам объекта. Среди получившихся пар отбрасываются те, чье расстояние превшает заданный порог, а число, равное количеству оставшихся пар, заносится в словарь с ключем, соответствующим идентификатору изображения.
        \item Словарь сортируется по убыванию значения, берутся первые r элементов. Идентифицируемое изображения относят к тому классу, к которому принадлженит большинство из данных r изображений.
        
    \end{enumerate}


\end{frame}

% --------------------
\begin{frame}{Эксперимент}
    В оригинальном эксперименте можно варьировать число r самых подходящих образцов, по которым определяется принадлежность к классу, а также порог, по которому принимаются соответствия между ключевым точками.
    Подбирая различные пороги, видно, что при низком пороге алгоритм плохо решает задачу, но с некоторого момента достигается высокая точкность. \\
    
    Гиперпараметр r не оказывает серьезного влияния на результат, при этом лучшие результаты достигаются при небольшом значении.
    \begin{columns}[T]
      \begin{column}{0.4\textwidth}
        \begin{figure}[H]
            \centering
            \includegraphics[scale=0.24]{acc_0.pdf}
            \caption{Зависимость точности от порога}
            \label{fig:orange}
         \end{figure}
      \end{column}
      \begin{column}{0.4\textwidth}
        \begin{figure}[H]
            \centering
            \includegraphics[scale=0.24]{acc_1.pdf}
            \caption{Зависимость точности от ранга r}
            \label{fig:orange}
         \end{figure}
      \end{column}
    \end{columns}

\end{frame}

% --------------------
\begin{frame}{Предлагаемая модификация}
    Несмотря на высокую точность, алгоритм имеет высокие затраты по времени.
    Идея улучшения состоит в том, что вместо сравнения ключевых точек у идентифицируемого изображения с каждым изображением из датасета, он сравнивается по точкам ветвления папиллярных линий, которых значительно меньше, и которые ищутся в скелете, представленном в виде графа. Итоговая точность составляет 95\%, что сравнимо с качеством без данной процедуры, Время на определение близости между двумя изображениями составляет : $876 * 10^{-5}$ секунды, в то время как у оригинального алгоритма: $11689 * 10^{-6}$, то есть модификация быстрее оригинала примерно в $\frac{4}{3}$ раза.  
    \begin{columns}[T]
      \begin{column}{0.4\textwidth}
        \begin{figure}[H]
            \centering
            \includegraphics[scale=0.05]{saved_image.png}
            \caption{Предобработанное изображение с выделенными ключевыми точками}
            \label{fig:orange}
         \end{figure}
      \end{column}
      \begin{column}{0.4\textwidth}
        \begin{figure}[H]
            \centering
            \includegraphics[scale=0.05]{skel_img.png}
            \caption{Скелет изображения с точками ветвления}
            \label{fig:orange}
         \end{figure}
      \end{column}
    \end{columns}
    

\end{frame}

% --------------------
\begin{frame}{Выводы}
    \begin{itemize}
        \item Показана эффективность предварительного отбора отпечатков
        \item Получено ускорение, без ущерба для точности
        
    \end{itemize}
    

\end{frame}

\begin{frame}{Конец}
    \begin{center}
        {\LARGE \textbf{Спасибо за внимание!}}
    \end{center}
\end{frame}



\end{document}